\section*{Chapitre \ref{ch:g4:breathing} --- La respiration}

\begin{solution}{exo:g4:breathing:1}
\omterm{def:g4:breathing:movements}{Inspiration}~: la cage thoracique se dilate et l'air entre.
\omterm{def:g4:breathing:movements}{Expiration}~: la cage thoracique se relâche et l'air sort.
\end{solution}

\begin{solution}{exo:g4:breathing:2}
Le nez (ou la bouche), la \omterm{def:g4:breathing:organs}{trachée}, les deux \omterm{def:g1:plants-around-us:tree}{branches}, les \omterm{def:g4:breathing:organs}{poumons},
et enfin les millions de minuscules poches d'air enveloppées de
vaisseaux sanguins.
\end{solution}

\begin{solution}{exo:g4:breathing:3}
Environ $15$ fois par minute. Dans un sprint, cela monte à $40$ ou
plus, et plus profondément --- les \omterm{def:g3:bones-and-muscles:muscle}{muscles} au travail ont besoin de
dioxygène plus vite, donc les \omterm{def:g4:breathing:organs}{poumons} en chargent plus vite.
\end{solution}

\begin{solution}{exo:g4:breathing:4}
Il prend du dioxygène et rend du dioxyde de carbone. L'échange se
fait dans les minuscules poches d'air des \omterm{def:g4:breathing:organs}{poumons}, à travers leurs
parois minces, avec le \omterm{prop:g4:journey-of-food:absorb}{sang}.
\end{solution}

\begin{solution}{exo:g4:breathing:5}
L'air expiré est plus chaud, plus humide (le miroir embué), plus
pauvre en dioxygène et plus riche en dioxyde de carbone.
\end{solution}

\begin{solution}{exo:g4:breathing:6}
Parce que l'échange air--sang a besoin de surface~: des millions de
parois de poches s'additionnent en un demi-court de tennis de
surface de rencontre repliée dans la poitrine. L'\omterm{def:g4:journey-of-food:tract}{intestin grêle}
emploie le même tour de force avec ses replis pour absorber les
\omterm{def:g4:journey-of-food:digestion}{nutriments}.
\end{solution}

\begin{solution}{exo:g4:breathing:7}
Le nez réchauffe, humidifie et filtre l'air en chemin --- un air
froid, sec ou poussiéreux arrive ainsi aux délicats \omterm{def:g4:breathing:organs}{poumons} déjà
préparé.
\end{solution}

\begin{solution}{exo:g4:breathing:8}
Réponse libre. En général environ $15$ au repos et $25$--$40$ après
les sauts~: les \omterm{def:g3:bones-and-muscles:muscle}{muscles} au travail ont réclamé plus de dioxygène,
donc la \omterm{def:g4:breathing:movements}{respiration} s'est accélérée et approfondie pour le fournir.
\end{solution}

\begin{solution}{exo:g4:breathing:9}
Que le corps stocke la nourriture généreusement (de quoi tenir des
heures ou des jours) mais ne garde presque aucune réserve de
dioxygène --- à peine une minute. Le dioxygène doit donc être livré
sans arrêt, souffle après souffle, ce qui explique que la
\omterm{def:g4:breathing:movements}{respiration} ne puisse jamais s'interrompre longtemps.
\end{solution}

\begin{solution}{exo:g4:breathing:10}
Vingt-cinq respirants ont pris du dioxygène à l'air de la salle et
lui ont rendu du dioxyde de carbone toute la matinée~: son air est
devenu plus pauvre en l'un et plus riche en l'autre. La
prescription est la règle 1~: ouvrir les fenêtres et aérer.
\end{solution}

\begin{solution}{exo:g4:breathing:11}
Les deux sont le lieu de rencontre du \omterm{prop:g4:journey-of-food:absorb}{sang} et du dioxygène~: de
vastes surfaces minces, tapissées de \omterm{prop:g4:journey-of-food:absorb}{sang}, où le dioxygène entre et
où le dioxyde de carbone sort --- les branchies prenant le dioxygène
dissous dans l'eau, les \omterm{def:g4:breathing:organs}{poumons} le prenant dans l'air. À l'air
libre, les surfaces plumeuses des branchies de la truite s'affaissent
et se collent~: l'immense surface de passage se réduit à presque
rien, et sans surface de passage, pas de dioxygène, si riche que
soit l'air autour.
\end{solution}
